% Part: proof-theory
% Chapter: normalization
% Section: introduction

\documentclass[../../../include/open-logic-section]{subfiles}

\begin{document}

\olfileid{pt}{nor}{seg}

\olsection{段与切割}

在推导中，一个引入规则之后紧接一个消去规则就构成一个\emph{绕道}，而将推导\emph{规范化}就是要消除这类绕道。
然而，\Elim{\lor} 和 \Elim{\lexists} 使得连“我们要消除的绕道”的定义也变得有些复杂。
问题在于：由于这些规则的存在，我们无法保证在一个绕道中，消去规则总是\emph{紧随}引入规则之后。
中间可能隔着一个 \Elim{\lor} 或 \Elim{\lexists}，例如：
\[
\AxiomC{}
\DeduceC{$\lexists[x][!C(x)]$}
\AxiomC{$\Discharge{!A}{x}$}
\DeduceC{$!B$}
\DischargeRule{\Intro{\lif}}{x}
\UnaryInfC{$!A \lif !B$}
\DischargeRule{\Elim{\lexists}}{y}
\BinaryInfC{$!A \lif !B$}
\AxiomC{}
\DeduceC{$!A$}
\RightLabel{\Elim{\lif}}
\BinaryInfC{$!B$}
\DisplayProof
\]
事实上，可以有任意多个 \Elim{\lor} 或 \Elim{\lexists} 将 \Intro{\lif} 与 \Elim{\lif} 隔开，例如：
\[
\AxiomC{}
\DeduceC{$\lexists[x][!C(x)]$}
\AxiomC{}
\DeduceC{$!D \lor !E$}
\AxiomC{$\Discharge{!A}{x}$}
\DeduceC{$!B$}
\DischargeRule{\Intro{\lif}}{x}
\UnaryInfC{$!A \lif !B$}
\AxiomC{$\Discharge{!A}{x}$}
\DeduceC{$!B$}
\DischargeRule{\Intro{\lif}}{x}
\UnaryInfC{$!A \lif !B$}
\RightLabel{\Elim{\lor}}
\TrinaryInfC{$!A \lif !B$}
\DischargeRule{\Elim{\lexists}}{y}
\BinaryInfC{$!A \lif !B$}
\AxiomC{}
\DeduceC{$!A$}
\RightLabel{\Elim{\lif}}
\BinaryInfC{$!B$}
\DisplayProof
\]
这个例子表明，同一个 \Elim{\lif} 可能参与多个绕道，甚至，因为它跟在不止一个~\Intro{\lif} 之后。
为了涵盖所有这些可能性，我们使用 \Log{N1i} 证明中出现的某类公式出现序列（称为\emph{段}）来定义绕道。
在我们的例子中，这就是一系列 $!A \lif !B$ 的出现，其中 \Elim{\lor} 或 \Elim{\lexists} 的小前提紧随其结论。
在该例中有两个这样的段，每个包含三个 $!A \lif !B$ 的出现。
以下是正式定义。

\begin{defn}
在 \Log{N1i} 推导~$\delta$ 中，一个\emph{段}是同一公式~$!A$ 在 $\delta$ 中的出现序列 $!A_1$, \dots, $!A_n$，满足以下条件：
\begin{enumerate}
  \item $!A_1$ 不是 \Elim{\lor} 或 \Elim{\lexists} 的结论；
  \item $!A_n$ 不是 \Elim{\lor} 或 \Elim{\lexists} 的小前提；
  \item 对 $i = 1$, \dots,~$n-1$，$!A_i$ 是 \Elim{\lor} 或 \Elim{\lexists} 的小前提，且 $!A_{i+1}$ 是该推理的结论。
\end{enumerate}
该段的\emph{长度}为~$n$。
\end{defn}

并非每个段都是要消除的目标（即绕道）。需要消除的段被称为\emph{切割}。

\begin{defn}
一个\emph{切割}是一个段，其中 $!A_n$ 是某个消去推理的主前提，并且要么 $n=1$ 且 $!A_1 = !A_n$ 是某个引入推理或~\FalseInt 的结论，要么~$n>1$。
\end{defn}

需要注意的是，一个切割段不一定以某个引入规则的结论开始。唯一的要求是它以某个消去规则的主前提结束。然而，一个长度为~1 的段必须本身是引入规则的结论，才能算作切割。

\begin{defn}
一个切割的\emph{秩}为~$\depth{!A}$。推导~$\delta$ 的\emph{切割秩}~$\cutrank{\delta}$ 是~$\delta$ 中所有切割的秩的最大值。一个切割在~$\delta$ 中是\emph{极大的}，如果其秩等于~$\cutrank{\delta}$，即其秩最大。$\delta$ 的\emph{切割长度}是其所有极大切割的长度之和。
\end{defn}

\end{document}